\section{Explainable AI}
\label{sec:xai}

\subsection{Scope and Design Choice}
Two complementary interpretability mechanisms are implemented in the system: native, model-intrinsic feature importance for the tree-based models, and a retrieval-augmented large-language-model layer that narrates the forecast in natural language. A design choice made explicit here is that the system does not compute post-hoc, model-agnostic attribution methods such as SHAP~\cite{lundberg2017shap} or LIME~\cite{ribeiro2016lime}; these are widely used in the broader explainable-AI literature (Section~\ref{sec:literature}) and are a natural direction for future extension of the feature-attribution component (Section~\ref{sec:future}), but were not implemented in the evaluated system, and this paper reports only interpretability mechanisms that are actually present in the codebase.

\subsection{Native Feature Importance}
For the gradient boosting and random forest models described in Section~\ref{subsec:ml-models}, the training procedure produces a gain-based (XGBoost) or impurity-based (random forest) importance score for every engineered feature as a direct byproduct of tree construction, at no additional computational cost. The system retains the ten most important features for the winning model, when the winning model is one of these two families, and surfaces them to the user as a ranked, labeled bar chart. Because these importances are computed identically to how the trees themselves were built, they are a \emph{global} explanation of which lag, rolling-statistic, calendar, or exogenous features the model relied upon in aggregate across the training data, rather than a \emph{local} explanation of any single forecasted value.

\subsection{Retrieval-Augmented Natural-Language Explanation}
The second, and more novel, interpretability mechanism is a retrieval-augmented generation layer built around a general-purpose large language model. Two independent context sources are combined before any natural-language response is generated. First, a set of \emph{deterministic} textual summaries are constructed directly and exclusively from the pipeline's own computed values: per-model accuracy and error metrics, the identity and score of the selected model, the forecast's minimum, maximum, and average values, a trend classification computed from a fixed percentage threshold, and, where applicable, per-product or per-group sales rankings and detected exogenous-feature usage. Second, for qualitative or advisory questions (for example, "how can I reduce stockouts"), the system retrieves the most semantically similar passages from a curated corpus of thirteen authored documents covering retail-operations topics such as inventory classification, safety stock, promotional and markdown strategy, seasonal demand planning, and stockout-versus-overstock diagnosis, using cosine similarity between text embeddings stored in a local, persistent vector index. This corresponds to the retrieval-augmented generation pattern of Lewis et al.~\cite{lewis2020rag}, in which language-model generation is conditioned on retrieved source text rather than relying solely on the model's parametric knowledge, reducing the risk of the response fabricating claims unrelated to the retrieved evidence.

\subsection{Global and Local Explanations}
In the terminology of the explainability literature, the tree-based feature-importance mechanism provides a global explanation of model behavior, while the natural-language layer can operate at both levels: a proactive business insights report narrates a global summary of the forecast and dataset, whereas the conversational assistant can answer local questions about a specific number, model, or product in the current result. Numeric questions, such as "which product sold the most," are answered exclusively from the deterministic, pre-computed per-product summary rather than from retrieval or from the language model's own estimation, since retrieval-based grounding for a fact the system already knows exactly would only introduce additional hallucination risk rather than reduce it; retrieval is reserved for genuinely qualitative, advisory content.

\subsection{Consistency Constraints and Trust}
A specific and, in the authors' assessment, important engineering measure is that the prompt supplied to the language model explicitly enumerates the system's own deterministic labels, most notably the trend classification and its fixed percentage threshold, and instructs the model never to re-derive or contradict them in its own words, even when the underlying numeric change is negative and could colloquially be described as a decline outside the labeled threshold band. The prompt additionally states plainly which models currently do and do not use exogenous features such as price, promotions, or holidays for the specific forecast being discussed, so that the language model cannot imply a capability the system did not actually exercise on this run. These constraints exist because a natural-language explanation that quietly contradicts the numbers the system itself computed would be more damaging to user trust than providing no explanation at all; grounding the language model in the pipeline's own deterministic output, rather than allowing it to reason independently over the raw numbers, is the mechanism by which this consistency is enforced in practice.

\subsection{How These Mechanisms Improve Trust}
Feature importance answers the question of \emph{which inputs the winning model relied upon}, while the grounded language layer answers the qualitatively different question of \emph{what the forecast means for the business and why a particular model was trusted}, including an explicit statement of confidence derived from the model's own cross-validated accuracy. Presenting both together, rather than either mechanism alone, allows a non-technical stakeholder to move from "the system predicts a number" to "the system predicts a number, driven primarily by these business factors, with this level of confidence, and here is what a comparable retail situation typically implies," which is the practical bar this system was designed to meet.
